\begin{table}[H]
\centering
\caption{Principal computational components used in the dissertation workflow.}
\label{tab:computational_environment}

\small
\renewcommand{\arraystretch}{1.20}
\setlength{\tabcolsep}{5pt}

\begin{tabular}{p{3.7cm}p{8.3cm}}
\toprule
\textbf{Component} & \textbf{Implementation} \\
\midrule

Execution environment
& Managed Google Colab notebook runtimes with GPU acceleration for model training and attribution generation and CPU-based execution for aggregation and statistical analysis \\

Persistent storage
& Google Drive dissertation project directories containing datasets, model artifacts, attribution maps, manifests, audit records, tables, and figures \\

Programming language
& Python \\

Deep-learning framework
& TensorFlow and Keras \\

Model architecture
& EfficientNet-B0 classifiers \\

Image processing
& OpenCV, NumPy, and Pillow \\

Data management
& Pandas and NumPy \\

Statistical analysis
& SciPy and Scikit-learn \\

Visualization
& Matplotlib and related Python plotting utilities \\

Attribution methods
& Grad-CAM, Grad-CAM++, Eigen-CAM, Score-CAM, Integrated Gradients, and SHAP \\

Artifact management
& Configuration-specific experiment directories, model checkpoints, prediction exports, numerical attribution maps, canonical manifests, remediation audits, and provenance-preserving dissertation exports \\

\bottomrule
\end{tabular}
\end{table}